\documentclass[10pt,a4paper]{article}
\usepackage[a4paper,margin=1cm,top=0.6cm,bottom=0.6cm]{geometry}
\usepackage{fontspec}
\usepackage[brazilian]{babel}
\usepackage{xcolor}
\usepackage{tikz}
\usetikzlibrary{calc, positioning, arrows.meta}
\usepackage{amsmath, amssymb}
\usepackage{booktabs}
\usepackage{multicol}
\usepackage{microtype}

% --- TIPOGRAFIA (Estilo Crisis) ---
\setmainfont{Helvetica Neue}[
  BoldFont={Helvetica Neue Bold},
  ItalicFont={Helvetica Neue Italic},
  Scale=0.9
]
\newfontfamily\headerfont{Helvetica Neue Condensed Bold}[Scale=1.4]
\newfontfamily\idfont{Helvetica Neue Condensed Bold}[Scale=3.5]
\newfontfamily\idsubfont{Helvetica Neue Condensed Bold}[Scale=1.6]
\newfontfamily\idlabelfont{Helvetica Neue Bold}[Scale=0.65]
\newfontfamily\sectionlabelfont{Helvetica Neue Bold}[Scale=0.45]
\newfontfamily\sectionnumfont{Helvetica Neue Condensed Bold}[Scale=0.9]

% --- CORES ---
\definecolor{crisisBlue}{RGB}{0, 150, 210} 
\definecolor{crisisOrange}{RGB}{230, 130, 40}
\definecolor{crisisGray}{RGB}{60, 60, 60}
\definecolor{crisisLightBlue}{RGB}{235, 245, 252}
\definecolor{crisisRed}{RGB}{180, 40, 40}
\definecolor{crisisLightRed}{RGB}{252, 235, 235}

\pagestyle{empty}
\setlength{\parindent}{0pt}

% --- MACROS DE LAYOUT ---

\newcommand{\crisisHeader}[5]{% #1=cor, #2=título, #3=subtítulo, #4=ID num, #5=ID letra
    \begin{tikzpicture}[remember picture, overlay]
        \fill[#1] ($(current page.north west)+(0,0)$) rectangle ($(current page.north east)+(0,-3.2cm)$);
        
        % Título
        \node[anchor=west, white] at ($(current page.north west)+(0.8cm,-1.6cm)$) {
            \begin{minipage}{12cm}
                {\fontsize{34}{38}\headerfont\MakeUppercase{#2}}\\[2pt]
                {\fontsize{18}{22}\headerfont\MakeUppercase{#3}}
            \end{minipage}
        };

        % Círculo ícone
        \draw[white, line width=1.5pt] ($(current page.north east)+(-5.5cm,-1.6cm)$) circle (0.8cm);
        \node[white] at ($(current page.north east)+(-5.5cm,-1.6cm)$) {
            \begin{tikzpicture}[scale=0.5, white]
                \draw[line width=1.5pt] (-0.3,0.3) arc (120:240:0.6) -- (0.3,-0.3) arc (-60:60:0.6) -- cycle;
                \draw[line width=1pt] (-0.1,0) -- (0.1,0); 
            \end{tikzpicture}
        };

        % ID BOX
        \node[fill=white, rounded corners=6pt, inner sep=6pt, anchor=north east] at ($(current page.north east)+(-0.4cm,-0.4cm)$) {
            \begin{minipage}{3.4cm}
                \centering
                \vspace{2pt}
                {\color{#1}\idfont #4\rlap{\raisebox{12pt}{\idsubfont #5}}}\\[-2pt]
                {\color{#1}\idlabelfont CRISIS PROTOCOLS}\\[2pt]
                {\setlength{\fboxsep}{2pt}\colorbox{#1}{\parbox{1.5cm}{\centering\color{white}\sectionlabelfont SECTION\\[-2pt]\sectionnumfont ONE}}%
                 \hspace{0.05cm}%
                 \colorbox{#1}{\parbox{1.5cm}{\centering\color{white}\sectionlabelfont SECTION\\[-2pt]\sectionnumfont TWO}}}\\[3pt]
                {\color{#1}\idlabelfont CRISIS PROCEDURES}\\[2pt]
            \end{minipage}
        };
    \end{tikzpicture}
    \vspace{3.0cm}
}

% Box de alerta estilo Crisis
\newcommand{\crisisAlert}[2]{% #1=título, #2=corpo
    \vspace{0.2cm}
    \begin{tikzpicture}
        \node[rectangle, draw=crisisRed, fill=crisisLightRed, rounded corners=2pt, inner sep=6pt, text width=\dimexpr\linewidth-16pt] {
            {\bfseries\small\textcolor{crisisRed}{#1}}\\[2pt]
            {\footnotesize #2}
        };
    \end{tikzpicture}
    \par\vspace{0.2cm}
}

% Box de sumário estilo Crisis
\newcommand{\crisisSummary}[2]{% #1=título, #2=corpo
    \vspace{0.2cm}
    \begin{tikzpicture}
        \node[rectangle, draw=crisisBlue, fill=crisisLightBlue, rounded corners=2pt, inner sep=6pt, text width=\dimexpr\linewidth-16pt] {
            {\bfseries\small\textcolor{crisisBlue}{#1}}\\[2pt]
            {\footnotesize #2}
        };
    \end{tikzpicture}
    \par\vspace{0.2cm}
}

\begin{document}

\crisisHeader{crisisBlue}{Oxímetro}{de Pulso SpO\textsubscript{2}}{42}{A}

\textbf{AUTOR:} Nick Mark MD | \textbf{TRADUÇÃO:} Equipe ICU\par
\vspace{0.3cm}

\setlength{\columnsep}{0.6cm}
\begin{multicols}{2}

\textcolor{crisisBlue}{\headerfont\large PRINCÍPIO E FISIOLOGIA}\par
\vspace{0.1cm}
{\small O oxímetro mede a saturação de oxigênio arterial ($SpO_2$) de forma contínua.}\par
\vspace{0.1cm}
{\small\bfseries \textcolor{crisisBlue}{1} Absorção Diferencial}: Utiliza luz \textcolor{red}{\textbf{Vermelha}} (660nm) e \textcolor{orange}{\textbf{Infravermelha}} (940nm).
\begin{itemize}
    \setlength{\itemsep}{0pt}
    \item {\small Hemoglobina Oxigenada ($HbO_2$) absorve mais Infravermelho.}
    \item {\small Hemoglobina Reduzida ($Hb$) absorve mais Vermelho.}
\end{itemize}
{\small\bfseries \textcolor{crisisBlue}{2} Fluxo Pulsátil (AC)}: Identifica o sangue arterial pelo pulso, filtrando o tecido/veias.

\vspace{0.2cm}
\begin{center}
\begin{tikzpicture}[scale=0.9, font=\sffamily\tiny]
    \fill[crisisLightBlue] (-0.3,-0.5) rectangle (6.2,3.2);
    \draw[->, crisisGray] (0,0) -- (5.5,0) node[right] {$\lambda$};
    \draw[->, crisisGray] (0,0) -- (0,3) node[above] {Abs.};
    \draw[thick, xcolor=red!80] (0.5, 2.5) .. controls (2, 2) and (4, 0.5) .. (5.0, 0.4) node[right] {$HbO_2$};
    \draw[thick, xcolor=blue!80] (0.5, 0.4) .. controls (2, 1) and (4, 2.5) .. (5.0, 2.5) node[right] {$Hb$};
    \draw[dashed, crisisGray!50] (1.1,0) -- (1.1,3); \node[above, scale=0.8] at (1.1,3) {660};
    \draw[dashed, crisisGray!50] (4.4,0) -- (4.4,3); \node[above, scale=0.8] at (4.4,3) {940};
\end{tikzpicture}
\end{center}

\crisisAlert{ALERTA DE SEGURANÇA: HIPOXEMIA OCULTA}{
    O oxímetro pode \textbf{superestimar} a saturação em pacientes com pele escura.
    \begin{itemize}
        \setlength{\itemsep}{0pt}
        \item Risco 3x maior de hipoxemia não detectada.
        \item SEMPRE correlacione com a clínica e gasometria ($SaO_2$).
    \end{itemize}
}

\textcolor{crisisBlue}{\headerfont\large HEMOGLOBINAS ANORMAIS}\par
\vspace{0.1cm}
\begin{itemize}
    \setlength{\itemsep}{0pt}
    \item {\small \textbf{Carboxihemoglobina ($COHb$)}: Falsamente interpretada como $HbO_2$. Valores normais mesmo em intoxicação severa por CO (incêndios).}
    \item {\small \textbf{Metemoglobina ($MetHb$)}: Leitura "travada" ao redor de \textbf{85\%}, ignorando o valor real.}
\end{itemize}

\vspace{0.3cm}
\textcolor{crisisBlue}{\headerfont\large ANÁLISE DA ONDA PLETISMOGRÁFICA}\par
\vspace{0.1cm}
{\small Valide o sinal observando a onda antes de agir sobre o valor númerico.}\par
\begin{itemize}
    \setlength{\itemsep}{0pt}
    \item {\small \textbf{Sístole}: Pico da onda. \textbf{Incisura Dicrótica}: Fechamento da valva aórtica.}
\end{itemize}
\vspace{0.1cm}
\begin{center}
\begin{tikzpicture}[scale=1.2, font=\sffamily\tiny]
    \draw[gray, opacity=0.2] (0,0) grid (5,1.2);
    \draw[thick, crisisBlue] (0,0.2) .. controls (0.2,0.2) and (0.3,1.1) .. (0.5,1.1) 
         .. controls (0.7,1.1) and (0.8,0.5) .. (0.9,0.6) 
         .. controls (1.0,0.7) and (1.4,0.2) .. (1.8,0.2)
         .. controls (2.0,0.2) and (2.1,1.1) .. (2.3,1.1)
         .. controls (2.5,1.1) and (2.6,0.5) .. (2.7,0.6) 
         .. controls (2.8,0.7) and (3.2,0.2) .. (3.6,0.2);
    \node[crisisBlue, font=\bfseries] at (0.5,1.3) {SÍSTOLE};
    \node[crisisGray] at (1.2,0.7) {Dicrótica};
\end{tikzpicture}
\end{center}

\crisisSummary{SUMÁRIO TÉCNICO \& REGRA DE OURO}{
    \begin{tabular}{@{}lp{5cm}@{}}
        \textbf{Choque} & Reduz sinal. Escolha nariz/orelha. \\
        \textbf{Malha} & Esmalte escuro interfere. \\
        \textbf{Lag Time} & Atraso de 15--60s periférico. \\
        \textbf{OURO} & \textbf{Gasometria Arterial}. \\
    \end{tabular}
}

\vspace{0.3cm}
\textcolor{crisisBlue}{\headerfont\large CURVA DE DISSOCIAÇÃO DA HEMOGLOBINA}\par
\vspace{0.1cm}
{\small Entenda a relação não-linear entre a $PaO_2$ (pressão dissolvida no plasma) e a $SpO_2$ (saturação).}\par
\begin{itemize}
    \setlength{\itemsep}{1pt}
    \item {\small \textbf{Platô (Segurança)}: $PaO_2$ > 60 mmHg mantém a saturação acima de 90\%. Oximetria é \textbf{insensível} a níveis perigosamente altos de oxigênio (ex: $PaO_2$ de 100 ou 500 mmHg resultarão ambos em $SpO_2$ $\sim$100\%).}
    \item {\small \textbf{Porção Descendente (Perigo)}: Quedas abaixo de 90\% ($PaO_2$ $<$ 60 mmHg) representam perda rápida e ameaçadora do conteúdo de oxigênio.}
    \item {\small \textbf{Desvio para a Direita}: Acidose, hipercapnia, hipertermia (facilita liberação para tecidos mas reduz afinidade no pulmão).}
    \item {\small \textbf{Desvio para a Esquerda}: Alcalose, hipotermia.}
\end{itemize}

\vspace{0.3cm}
\textcolor{crisisBlue}{\headerfont\large LIMITAÇÕES NA PREVALÊNCIA DE FALÊNCIA RESPIRATÓRIA}\par
\vspace{0.1cm}
{\small Em pacientes com SDRA (Síndrome do Desconforto Respiratório Agudo) ou pneumonia viral grave:}
\begin{itemize}
    \setlength{\itemsep}{1pt}
    \item {\small \textbf{Hipoxemia Permissiva}: Alguns protocolos sugerem alvo de 88-92\% para evitar toxicidade pelo oxigênio e injúria pulmonar associada ao ventilador (VILI), desde que não haja sinais de choque, isquemia miocárdica ou disfunção cognitiva aguda.}
    \item {\small \textbf{Respiração Paradoxal (Silent Hypoxia)}: Pacientes podem apresentar $SpO_2$ profundamente baixas (<70\%) com pouca dispneia. Monitorar o trabalho respiratório e FR é tão importante quanto o valor da oximetria numérico isolado.}
\end{itemize}

\end{multicols}

\vfill
\begin{center}
    {\fontsize{7}{9}\selectfont \textcolor{crisisGray}{© 2025 ICU Crisis Protocols - Adaptado de Nick Mark MD. Uso educacional.}}
\end{center}

\end{document}
